\section{范围：如何改进研究过程}
\label{sec:scope}

本文关心如何构建越来越会做研究的 AI：选择问题，把直觉转成可检验的命题，获得证据，修订判断，并改善后续决策。有用的系统不必自动化每个环节。帮助研究者放弃一次没有信息量的实验，可能比生成一篇完整但薄弱的论文更有价值。

我们用 \emph{autoresearch} 指这一广义方法方向，其中包括人机协作和自主系统。\emph{Meta-science} 研究科学如何运行，以及如何改进科学实践。机器学习中较窄意义的 \emph{meta-learning}，则是跨任务学习科研策略的一条技术路线，并不是所有收录方法都必须采用的机制。本文以研究过程为对象，以 AI research 为主要实验场景。数学提供补充 testbed，因为其中很多候选答案、程序和构造可以被精确检查。可检查性有助于评价，但不能决定问题本身是否重要。

\subsection{五个问题与两个时间尺度}

\begin{figure}[t]
\centering
\begin{tikzpicture}[>=Latex,box/.style={draw=inkblue!60,fill=softblue,rounded corners=2pt,align=center,text width=3.05cm,minimum height=1.15cm,font=\small\sffamily},every path/.style={line width=0.6pt}]
\node[box] (taste) {Q1：科研品味\\选择值得研究的问题};
\node[box,right=0.45cm of taste] (question) {Q2：问题形成\\形成可检验的命题};
\node[box,right=0.45cm of question] (execute) {Q3：研究执行\\获得可靠证据};
\node[box,right=0.45cm of execute] (learn) {Q4：经验学习\\改善后续决策};
\draw[->] (taste)--(question);
\draw[->] (question)--(execute);
\draw[->] (execute)--(learn);
\draw[->] (learn.north)--++(0,0.5)-| (taste.north);
\node[draw=green!40!black,fill=softgreen,align=center,text width=14cm,below=0.65cm of $(question.south)!0.5!(execute.south)$,font=\small,minimum height=0.7cm] (eval) {Q5：贯穿整个流程的评价：价值、有效性、成本与迁移};
\draw[dashed,->] (eval.north)--($(question.south)!0.5!(execute.south)$);
\end{tikzpicture}
\caption{本文的研究地图。反馈既可以修订当前项目，也可以改变后续项目的研究方式。顺序用于组织讨论，不是固定工作流。}
\label{fig:map}
\end{figure}

Q1 讨论哪些问题值得投入，Q2 讨论怎样形成能通过研究解决的问题，Q3 讨论系统能否产生支持结论的证据，Q4 讨论经验究竟改变了研究者的什么，Q5 则讨论如何分别测量这些能力。完整 AI scientist systems 放在 Q3 中讨论，因为它们展示了各环节如何相互影响。

这里存在两个不同的时间尺度。在一个项目内部，agent 可以改进候选解、修订实验或切换搜索分支。在不同项目之间，它可以保留报告、积累技能、修改控制代码或更新模型权重。评价后一类改善时，可以采用下面的概念性比较：
\begin{equation}
\Delta_{\mathrm{transfer}} = S(\pi_{\mathrm{after}};\mathcal{T}_{\mathrm{new}},B)-S(\pi_{\mathrm{before}};\mathcal{T}_{\mathrm{new}},B),
\end{equation}
其中 $\mathcal{T}_{\mathrm{new}}$ 包含留出的任务家族，$B$ 为相同的评价预算，$S$ 是事先声明的结果指标。这是本文建议的比较方式，不是科学进步的通用定义；获取经验本身的成本也需要报告。在适应所用的同一任务上获得更好结果，不能估计这一迁移量。

\subsection{什么是好的研究问题？}

问题可能新颖却不重要，重要却暂时无法检验，也可能容易执行但已经有答案。因此我们分别讨论五个判断：解决未知的\textbf{重要性}，相对于相关工作的\textbf{新颖性}，命题的\textbf{可检验性}，给定预算下的\textbf{可行性}，以及包括负结果在内的\textbf{信息价值}。它们是本文的分析维度，并非已经验证的评分量表；各维度的权重取决于领域和研究目的。

例如，“提高这个分类器的准确率”给出了优化目标。“一种新 augmentation 的收益来自不变性，还是仅仅来自更多训练更新？”则识别了相互竞争的解释。后者自然引出 matched-compute control，即使新方法没有增益，也可能产生有用知识。这个例子用于说明区别，不表示该 augmentation 问题具有新颖性。实际挑战在于让 agent 根据真实文献和可获得证据，提出这样的区分。

研究者的文章能帮助描述目前 benchmarks 只测到一部分的判断力。Olah 提出用低成本、带噪声的反馈训练 taste；Karpathy 讨论项目选择与推进；Hamming 强调重要问题的选择~\cite{colah-taste-exercises,karpathy-phd-guide,hamming-you-and-your-research}。Nielsen 强调好问题经常需要在工作中被发现，\emph{Research Debt} 则说明更好的抽象与解释也可以成为研究贡献~\cite{nielsen-research-heuristics,research-debt}。Nielsen 与 Qiu 进一步把科学判断放在注意力和资源分配制度中讨论~\cite{vision-of-metascience}。这些观点可以启发 AI 研究假设，但不是已经证明能训练有效 agent 的实验证据。

\subsection{文献范围与综合方式}

本文是基于固定仓库快照的批判性叙述综述，并不声称完成了穷尽性的 systematic review。快照包含 72 篇论文、6 篇文章或演讲，以及 Karpathy 的 autoresearch 软件项目。表~\ref{tab:corpus} 按仓库中的主要贡献分类计数。79 项资源全部进入讨论与参考文献；有些用于背景理解，而非作为核心实证证据。配套源码保存条目标识、主要来源链接、核查日期和引用到章节的覆盖关系。

\begin{table}[t]
\centering\small
\caption{按仓库主类别统计的文献集合。类别在概念上可以重叠，每项资源只计数一次。}
\label{tab:corpus}
\begin{tabularx}{\linewidth}{Y r Y r}
\toprule
类别 & 数量 & 类别 & 数量\\
\midrule
科研品味 & 8 & 问题形成 & 14\\
研究执行 & 11 & 经验学习 & 18\\
评价 & 16 & 完整系统 & 6\\
文章与演讲 & 6 & 总计 & 79\\
\bottomrule
\end{tabularx}
\end{table}

比较方法时，我们关注 agent 负责哪些决策、人类提供哪些输入、系统使用什么反馈，以及研究如何评价成功。关键方法论描述通过原始论文或官方材料核对。来源核查不等于独立实验复现。不同 backbone、预算、任务版本和评审协议下的 headline scores 不直接汇总成排行榜。每个问题章节末尾的方案是本文提出的研究设计，其中样本数和资源上限属于建议的 pilot 设置，不是统计功效计算或已经证明最优的配置。正式声称方法新颖性之前，仍需围绕具体方案补做检索。
